\documentclass[11pt]{article}
\usepackage[margin=1in]{geometry}
\usepackage{amsmath,amssymb,amsthm,hyperref}
\newtheorem{theorem}{Theorem}
\newcommand{\E}{\mathbb E}
\title{A quantitative stability gap for the residue-one collision bound}
\author{Joshua Mark}
\date{20 November 2026; research draft, priority not established}
\begin{document}\maketitle
Let $n\ge5$, $n\equiv1\pmod4$, $N=n+3$, and
$A=(n-1)(n+3)=N(N-4)$. Let $w$ be a probability distribution on
$\{\pm1\}^{n-1}$, with distinct sign vectors regarded as distinct atoms.
Write $v_x=(1,x)$ and assume exact isotropy
$\sum_x w_xv_xv_x^T=I_n$. Define
\[
 C=\sum_xw_x^2,\quad \delta=1/N-C,\quad
 \epsilon_x=\prod_jx_j,\quad z=\sum_xw_x\epsilon_xv_x,
 \quad h_x=\epsilon_x-v_x\cdot z.
\]
Only positive-weight atoms are used below.

\begin{theorem}[Nonnegative defect decomposition]
Put $q_{xy}=\epsilon_x\epsilon_y(v_x\cdot v_y)$ and
$f(q)=q^2+2q-3$. Then
\begin{equation}\label{defect}
 A\delta=2\sum_xw_xh_x^2+
 \sum_{x\ne y}w_xw_y f(q_{xy}).
\end{equation}
Both terms are nonnegative. In particular $\delta\ge0$, and the
ordered mass of off-diagonal pairs with $q_{xy}\notin\{-3,1\}$ is
at most $A\delta/32$. Furthermore
\begin{equation}\label{disp}
 \sum_xw_x|w_x-1/N|\le \sqrt{2\delta/A}+\delta.
\end{equation}
\end{theorem}
\begin{proof}
If $x,y$ differ at $r$ coordinates, then
$q_{xy}=(-1)^r(n-2r)\equiv1\pmod4$. On these integers $f\ge0$;
it vanishes precisely at $-3,1$, and otherwise is at least $32$.
Isotropy gives $\E q^2=n$ and $\E q=\|z\|^2$.
Bessel's inequality for the orthonormal functions comprising $v$ gives
$\|z\|^2\le1$, and direct expansion gives
$\E h^2=1-\|z\|^2$. As $f(q_{xx})=A$, separating the diagonal in
$\E f(q)=n-3+2\|z\|^2$ proves \eqref{defect}.
For a fixed support atom $x$, set
$B_x=\sum_{y\ne x}w_yf(q_{xy})\ge0$. Its conditional expectation is
\[
 Aw_x+B_x=n-3+2\epsilon_x v_x\cdot z
 =n-1-2\epsilon_xh_x.
\]
Since $A/N=n-1$, this gives
$A(w_x-1/N)=-2\epsilon_xh_x-B_x$. Average absolute values,
apply Cauchy--Schwarz, and use \eqref{defect} to obtain
\[
 \sum w|w-1/N|\le (2\sqrt{\E h^2}+\E B)/A
 \le\sqrt{2\delta/A}+\delta.
\]
\end{proof}

\begin{theorem}[Explicit extraction from a sufficiently small deficit]
If $\delta\le1/(1024N^4)$, the set
\[
 G=\{x:|Nw_x-1|\le1/(4N)\}
\]
has exactly $N$ members. Giving those members equal weight $1/N$
preserves exact isotropy. Their augmented sign rows can consequently
be completed by three sign columns to a Hadamard matrix of order $N$.
Thus, if no Hadamard matrix of order $N$ exists, every such isotropic
distribution has $C<1/N-1/(1024N^4)$.
\end{theorem}
\begin{proof}
Set $e=\sum_xw_x|Nw_x-1|$. Since $A\ge N^2/2$ and $N\ge8$,
\eqref{disp} implies
\[
 e\le N\delta+N\sqrt{2\delta/A}
 \le\frac1{1024N^3}+\frac1{16N^2}<\frac1{8N^2}.
\]
With $\eta=1/(4N)$, Markov's inequality bounds the mass $b$ outside
$G$ by $e/\eta<1/(2N)$. Each member of $G$ has weight between
$(1-\eta)/N$ and $(1+\eta)/N$. There cannot be $N+1$ such members,
since $(N+1)(1-\eta)/N>1$. There cannot be at most $N-1$, since
their mass would be at most $(N-1)(1+\eta)/N<1-1/(2N)$.
Hence $|G|=N$.
The $\ell^1$ distance from $w$ to the uniform distribution on $G$ is
at most $b+\eta<3/(4N)$. Every matrix entry in the resulting isotropy
error is bounded by this distance, since $|(v_x)_i(v_x)_j|=1$.
But $N$ times that error is an integer, because the uniform Gram
entries are integers divided by $N$ and the target is $I_n$.
Its absolute value is below one, so it is zero.

For completeness, apply the zero-defect case of \eqref{defect} to
this uniform isotropic distribution. Its switched rows
$V'_x=\epsilon_xv_x$ have off-diagonal inner products $-3$ or $1$.
The residual $K=NI_N-V'V'^T$ is positive semidefinite of rank three,
with diagonal $3$ and off-diagonal $3$ or $-1$. Real Gram vectors for
$K$ with inner product $3$ coincide; different types have Gram
$4I-J$, hence there are at most four types. Assign the four sign
triples $(1,1,1),(1,-1,-1),(-1,1,-1),(-1,-1,1)$ to these types.
Their Gram is $K$, so appending their coordinates to $V'$ produces
the claimed Hadamard matrix. This is an explicit extraction, not an
assumption that deleting and renormalizing preserves isotropy.
\end{proof}

The theorem requires an exactly isotropic input distribution with the
stated near-extremal collision value. No such input is constructed in
unknown Hadamard orders. The numerical constant is deliberately
conservative. The argument is elementary and its priority is not
established; it should be compared with weighted orthogonal-array and
linear-programming stability literature before any novelty claim.
\end{document}
